% Generated draft: requires mathematical and source review.
\documentclass[11pt]{article}
\usepackage[margin=25mm]{geometry}
\usepackage{fontspec}
\setmainfont{Latin Modern Roman}
\usepackage{amsmath,amssymb,mathtools,unicode-math}
\setmathfont{Latin Modern Math}
\usepackage{enumitem,needspace,xurl,hyperref}
\setlength{\parindent}{0pt}
\setlength{\parskip}{5pt}
\setlength{\emergencystretch}{3em}
\newcommand{\N}{\mathbb{N}}
\newcommand{\Z}{\mathbb{Z}}
\newcommand{\Q}{\mathbb{Q}}
\newcommand{\R}{\mathbb{R}}
\newcommand{\C}{\mathbb{C}}
\newcommand{\F}{\mathbb{F}}
\newcommand{\A}{\mathbb{A}}
\newcommand{\PP}{\mathbb{P}}
\newcommand{\E}{\mathbb{E}}
\newcommand{\eps}{\varepsilon}
\newcommand{\dd}{\,\mathrm{d}}
\begin{document}
\section*{UnsolvedMath ID: 6700074}
\textbf{Source problem code:} catalogue entry
\subsection{Definitions and mathematical statement}
Let \((W,g)\) be a smooth Riemannian manifold of dimension \(n+1\), and measure Hausdorff dimension and Hausdorff measure with its Riemannian distance. The conjecture concerns every closed subset \(Y\subset W\) with
\[
 \dim_{\mathrm H}(Y)<n-1=\dim(W)-2.
\]
It asserts that there are decreasing smooth domains \(V_1\supseteq V_2\supseteq\cdots\) containing \(Y\), with
\[
 Y=\bigcap_{i=1}^{\infty}V_i,\qquad
 \operatorname{Scal}_{g|_{\partial V_i}}\longrightarrow+\infty.
\]
The natural uniform meaning of the last condition is: for every \(A\in\R\), for all sufficiently large \(i\), every \(p\in\partial V_i\) satisfies \(\operatorname{Scal}_{g|_{\partial V_i}}(p)>A\). This is the intrinsic scalar curvature of the induced boundary metric, not boundary mean curvature. Uniformity is not explicitly defined in the excerpt.

The source says the claim is obvious for \(n=2\), and separately reports the same representation when \(Y\) lies in a smooth hypersurface \(Z^n\subset W\) and \(\mathcal H^{n-1}(Y)=0\). These are reported source assertions, not proved here.

The passage then starts defining \(\operatorname{Sc}_g(\partial\downarrow Y)\), but ends before the curvature inequality. Its likely intended completion is
\[
 \sup\{\sigma\in\R:\ \forall\text{ neighborhoods }U\supset Y,\ \exists\text{ smooth domain }V,
 \ Y\subset V\subset U,\ \operatorname{Scal}_{g|_{\partial V}}\geq\sigma\text{ everywhere}\}.
\]
The final inequality is a marked reconstruction of a truncated sentence. The dimension range, compactness or connectedness requirements on domains, and treatment of noncompact boundaries are not supplied; the definition should not be certified without the missing text.
\subsection{Short English statement}
Can every closed subset of Hausdorff dimension below ambient dimension minus two be obtained as a decreasing intersection of smooth domains whose induced boundary scalar curvatures tend uniformly to infinity? The excerpt also begins, but does not finish, a related local curvature invariant.
\subsection{Sources}
\begin{itemize}
\item Local UnsolvedMath registry, permanent ID 6700074, source code catalogue entry; Misha Gromov, 2017; updated 2026-08-06.
\item Gromov, \emph{101 Questions, Problems and Conjectures around Scalar Curvature} (2017), printed tag [?78], occurrence 77, PDF page 78, ``Subsets with Low Hausdorff Dimensions are Remains of Scalar Curvature Blow-ups'': \url{https://www.ihes.fr/~gromov/wp-content/uploads/2018/08/101-problemsOct1-2017.pdf}. Registry access date 2026-07-29; normalized extraction from an incomplete public author-hosted PDF; present status and reuse terms need review.
\item The supplied open-status summary mentions the claimed \(n=2\) case and particular constructions but supplies no further exact citation.
\end{itemize}
\end{document}
